\section{Conclusion}

\label{sec:conclusion}

We have identified temporal dominance as a structural failure mode in the deployed hybrid memory selector, derived the exact algebraic condition under which it occurs, and proposed CTLS as a principled correction. The deployed hybrid is outperformed by BM25 at all three evaluated token budgets in a controlled evaluation where only the selector changes and all other variables are held fixed. The temporal term causes most of this deficit: removing it (hybrid\_no\_time) recovers the majority of the gap. CTLS eliminates temporal dominance by construction through tier partition, preserving temporal decay as a within-tier tiebreaker rather than discarding it globally.

The Non-Dominance Property proved in Section~\ref{sec:method_formal} guarantees that no zero-overlap item can outrank a positive-overlap item under CTLS, regardless of freshness gap, frequency, or importance values. This is a structural guarantee, not a hyperparameter constraint: it holds for any positive value of the temporal weight within the relevant-tier scoring function. CTLS requires no model calls, no embedding model, no learned parameters, and no write-path changes. It is a zero-cost selector substitution that can replace the deployed hybrid in any context that uses the same frozen-bank, whole-item-packing retrieval architecture.

\subsection{Broader Research Directions}

\label{sec:conc_future}

Several research directions follow from this work. The most immediate is a direct implementation and evaluation of CTLS on the same frozen bank under the same controlled conditions: generating CTLS-ranked candidate lists, running the same answering model under the same token budgets, and computing cluster bootstrap intervals for the CTLS-versus-BM25 comparison. The present artifact contains all the inputs needed for that experiment: the frozen bank, segment hashes, metadata, token counts, and question-level records.

The adaptive query-conditioned weight formulation described in the method section is a natural extension of CTLS. Rather than a hard binary partition, a smooth version of the weight schedule could vary $w_t$ continuously with the mean lexical overlap of the candidate set. Evaluating whether this soft variant outperforms the hard partition -- especially on queries where most candidates have small but nonzero Jaccard overlap -- would establish whether the binary cut is the right design choice or merely the simplest one.

Multi-session and temporal-update question types in LongMemEval present a different operating point for recency versus relevance. In these settings, a record from a more recent session may genuinely contain the correct updated fact, making freshness a useful proxy for relevance rather than a noise source. Extending the mechanism analysis to these cases requires separating the causal structure of the question -- whether the answer is in the most recent relevant session or the oldest -- from the global freshness bias that the deployed hybrid introduces uniformly. A tier-partitioned selector could handle both cases cleanly if the partition criteria were extended beyond lexical overlap to include temporal proximity to a detected reference date in the question.

Finally, the connection between CTLS and the binary knapsack formulation of budget-constrained packing deserves exploration. CTLS changes the order in which items are offered to the packer; combining it with an exact packing policy that maximizes total relevance score within the budget constraint could close the remaining gap between CTLS and the theoretical optimum. With at most ten candidates per query, the exact knapsack solution is tractable at inference time and adds no model calls.
 The diagnosis identifies the exact condition under which a common retrieval heuristic fails, explains why the failure is invisible when evaluated only against a weak baseline, and shows that a large margin against a strong lexical baseline plus a clean ablation together pinpoint the cause. The fix is provably correct, immediately actionable, and requires no new infrastructure. Future work should implement CTLS directly and evaluate it on the existing frozen bank under the same controlled conditions, complete the blinded judge adjudication, extend the mechanism analysis to multi-session and temporal-reasoning question types where the dynamics of recency and relevance may differ, and measure the lifecycle cost benefit of improved selection quality across the full write-amortization curve.

